% Cấp độ 3: MMKG tiếng Việt tự xây và chuyển giao phương pháp cấp 2.
% Nguồn số liệu: vietnamese-mmkg/DATA_CARD.md, experiments/data_stats.md (khối vnmmkg),
% experiments/checks/multilingual_clip_probe.json, experiments/level2/val/choice.json,
% experiments/level3/vnmmkg/xlmr_h14/{a0_metrics,a1_metrics,a0_explanations}.json.
\section{MMKG tiếng Việt tự xây và phân tích chuyển giao (cấp độ 3)}\label{sec:vnmmkg}

Cấp độ 3 yêu cầu tự xây một MMKG (đồ thị tri thức có thực thể kèm ảnh) nhỏ cho một miền hẹp, ưu tiên tiếng Việt,
rồi chạy phương pháp tốt nhất của cấp 1--2 trên đó. Data card (nguồn, cách tạo triple, giấy phép, điều kiện sử
dụng) ở \file{vietnamese-mmkg/DATA_CARD.md}.

\subsection{Dữ liệu}\label{sec:vn-data}

\textbf{Miền và nguồn.} Di tích và công trình ở Việt Nam (chùa, đền, tháp, nhà thờ, lăng, bảo tàng, nơi có danh
hiệu di sản), vì miền này có dữ liệu mở và nhiều quan hệ không thấy được trong ảnh: ở tỉnh nào, xây thời nào, ai
sáng lập, thờ ai. Thực thể và triple lấy từ Wikidata~\cite{vn-vrandecic2014wikidata} (CC0); phần mở đầu và
infobox lấy từ Wikipedia tiếng Việt~\cite{vn-viwiki} (CC BY-SA 4.0); ảnh lấy từ Wikimedia
Commons~\cite{vn-commons}, tải ngày 24--25/10/2026.

\textbf{Quy mô} (khung gợi ý của đề trong ngoặc): 150 thực thể (50--150), gồm 75 di tích và 75 thực thể ngữ cảnh
(tỉnh, người, tôn giáo, danh hiệu, thời kỳ, lớp,~\ldots); 325 triple (300--1\,000) thuộc 13 quan hệ; 331 ảnh
(200--500), mỗi di tích 3--5 ảnh. Di tích là thực thể Wikidata ở Việt Nam có danh hiệu di sản hoặc thuộc một
trong 10 lớp công trình, xếp theo số sitelink.

\textbf{Triple.} Mỗi bộ ba giữ một nguồn, ưu tiên Wikidata, rồi infobox, rồi suy diễn. 257 triple là giá trị
truthy (hạng tốt nhất) trên Wikidata; 26 triple đọc từ ba trường infobox (người sáng lập, thờ, tôn giáo), chỉ nhận
giá trị là wikilink để ánh xạ được sang QID; 42 triple sinh từ hai quy tắc cố định: tỉnh (theo chuỗi
\texttt{P131}, chỉ tạo khi tới đúng một tỉnh) và thời kỳ (năm dựng đối chiếu với bảng 8 mốc triều đại). Toàn bộ
107 triple từ infobox và suy diễn được duyệt tay; 1 triple bị loại.

\textbf{Liên kết ảnh -- thực thể.} Ứng viên là ảnh \texttt{P18} cộng các file nằm trực tiếp trong Commons category
của di tích (JPEG/PNG rộng từ 400\,px, tối đa 12). Từng ảnh được duyệt tay, chỉ giữ ảnh thể hiện chính di tích;
bản đồ, ảnh gần trùng, ảnh có người là chủ thể chính bị bỏ (541 ảnh). Ảnh dùng chung chỉ giữ ở di tích xếp trên,
nên mỗi ảnh thuộc đúng một di tích; 35/110 di tích bị loại vì còn dưới 3 ảnh.

\textbf{Truy vấn.} 197 câu \emph{tự động} là tối đa 3 câu đầu bài Wikipedia, đã che nhãn và alias. 30 câu
\emph{soạn tay} gồm 12 câu trực tiếp (tả điều thấy trong ảnh) và 18 câu gián tiếp (cần ít nhất một quan hệ không
thấy trong ảnh); đáp án câu gián tiếp được tính tự động bằng cách khớp mẫu quan hệ trên triple. Câu đúng ở hạng
$k$ nếu trong $k$ ảnh đầu có ảnh của một di tích đáp án.

\textbf{Giấy phép và kiểm tra.} Chỉ nhận ảnh CC0, Public Domain, CC BY, CC BY-SA (163 ảnh CC BY-SA, 83 CC BY, 57
Public Domain, 28 CC0); tác giả, giấy phép và URL nguồn từng ảnh ghi trong \file{images.jsonl}. Bộ dữ liệu chỉ
dùng cho học thuật trong môn học, phải ghi công tác giả khi hiển thị ảnh. Quyết định duyệt ảnh và triple lưu
trong \file{review_images.json}, \file{review_triples.json}; chạy lại \texttt{python -m src.data.build\_vnmmkg}
từ cache cho ra đúng các file đã commit.

\subsection{Chống rò rỉ và giao thức đánh giá}\label{sec:vn-leak}

Bộ này quá nhỏ để vừa chọn tham số vừa chấm điểm, nên nhóm chốt hai giao thức trước khi chạy:
\begin{itemize}
\item \textbf{A0, zero-shot (chính):} không chọn gì trên bộ tiếng Việt. Pool 331 ảnh, 197 câu tự động và 30 câu
  soạn tay. Encoder chọn trên MKG-W val (mục~\ref{sec:vn-encoder}); $\alpha=0$ lấy từ
  \file{experiments/level2/val/choice.json} (hàng \texttt{kg\_triple}, cả ba seed).
\item \textbf{A1, chọn lại $\alpha$ (phụ):} chia di tích 50/50 bằng seed 0 (val: 38 di tích, 170 ảnh, 94 câu; test:
  37 di tích, 161 ảnh, 103 câu). $\alpha$ quét từ 0 đến 1, bước 0{,}05, chọn theo R@1 trên val (hoà thì theo
  MRR); test chạy một lần. Câu soạn tay không tham gia.
\item Câu soạn tay chỉ dùng để chấm: \file{queries_hand.json} khoá ở commit \texttt{55fa697} (26/10/2026), trước
  mọi lần chạy, không sửa sau đó.
\item Triple là tri thức có sẵn, không phải cạnh thống kê đồng xuất hiện, nên được dùng hết.
\item Encoder đóng băng, không huấn luyện: kết quả tất định, chạy một lần, không có $\pm$std. Độ bất định đến từ cỡ
  mẫu: một câu bằng 5{,}56 điểm ở nhóm gián tiếp, 8{,}33 điểm ở nhóm trực tiếp, 0{,}97 điểm ở test A1.
\end{itemize}

\subsection{Chọn encoder cho truy vấn tiếng Việt}\label{sec:vn-encoder}

CLIP ViT-B/32 của OpenAI (encoder cấp 1--2) chỉ học văn bản tiếng Anh. Nhóm so bốn encoder trên 84 thực thể MKG-W
val có mô tả tiếng Việt, mỗi thực thể hỏi bằng cả tiếng Việt và tiếng Anh, pool 892 ảnh
(Bảng~\ref{tab:vn-encoder}). Encoder đa ngữ gồm XLM-R~\cite{vn-conneau2020xlmr} ghép tháp ảnh huấn luyện trên
LAION-5B~\cite{vn-schuhmann2022laion5b} (bản OpenCLIP~\cite{vn-ilharco2021openclip}) và một bộ mã hoá văn bản đa
ngữ chưng cất (knowledge distillation) vào không gian CLIP B/32~\cite{vn-reimers2020multilingual}.

\begin{table}[t]
\centering
\caption{Chọn encoder trên MKG-W val (84 thực thể, pool 892 ảnh). Tất định, chạy một lần. MRR nhân 100.}
\label{tab:vn-encoder}
\footnotesize
\setlength{\tabcolsep}{3pt}
\begin{tabular}{@{}lrrrrr@{}}
\toprule
 & \multicolumn{3}{c}{Tiếng Việt} & \multicolumn{2}{c}{Tiếng Anh} \\
\cmidrule(lr){2-4}\cmidrule(l){5-6}
Encoder (văn bản + ảnh) & R@1 & R@10 & MRR & R@1 & MRR \\
\midrule
OpenAI ViT-B/32 (cấp 1--2) & 10,71 & 17,86 & 14,10 & 35,71 & 46,03 \\
XLM-R base + ViT-B/32 & 25,00 & 61,90 & 37,83 & 28,57 & 40,50 \\
Chưng cất đa ngữ + ViT-B/32 & 28,57 & 55,95 & 40,05 & 33,33 & 44,62 \\
XLM-R large + ViT-H/14 & \textbf{40,48} & \textbf{78,57} & \textbf{51,80} & 40,48 & 51,31 \\
\bottomrule
\end{tabular}
\end{table}

Với câu tiếng Việt, R@1 của OpenAI B/32 giảm từ 35,71 xuống 10,71, nên không dùng được. XLM-R large + ViT-H/14
cao nhất và không chênh giữa hai ngôn ngữ; nhóm chốt encoder này. Hai encoder B/32 đa ngữ chỉ cách nhau 3 câu
(mức nhiễu), còn H/14 hơn chúng khoảng 12 điểm R@1 (10 câu). Cấp 1--2 vẫn giữ OpenAI B/32, nên cột ``CLIP'' ở mục
này không so trực tiếp được với mục~\ref{sec:l2}.

\subsection{Phương pháp chuyển giao}\label{sec:vn-method}

Phương pháp được chuyển là kênh đếm cạnh KG (hàng \texttt{kg\_triple}, mục~\ref{sec:method-l2}), mô hình tốt
nhất ở cấp 2: trên MKG-W test nó nâng R@1 từ 24,89 lên 29,82 (mục~\ref{sec:l2}). Bộ xếp hạng GAT của cấp 1 cần
scene graph của từng ảnh, còn bộ này chỉ có KG mức thực thể. Kênh embedding NativE phải train lại trên 325
triple, quá ít để học embedding; hơn nữa trên MKG-W val kênh này nhận trọng số $\beta=0$, nên mô hình chính của
cấp 2 trùng với \texttt{kg\_triple}.

Chỉ thay đồ thị, ngôn ngữ và encoder. Linker (\texttt{EntityLinker}) chuyển câu về chữ thường, tách theo chữ cái
(mỗi âm tiết là một token, giữ dấu), so khớp chính xác với nhãn và alias tiếng Việt; độ dài tên tối đa nâng từ 5
token lên 12 âm tiết. Điểm đồ thị của một ảnh trong top-50 là số cạnh KG trực tiếp giữa di tích của ảnh và các
thực thể được nhắc; điểm cuối là $\alpha\, z(\mathrm{CLIP}) + (1-\alpha)\, z(\text{đồ thị})$. Với $\alpha=0$,
thứ tự CLIP chỉ còn dùng để phá thế hoà. Giải thích là chính danh sách cạnh đã đếm. Code:
\file{scripts/eval_vnmmkg_transfer.py}.

\subsection{Kết quả}\label{sec:vn-results}

\begin{table}[t]
\centering
\caption{A0, zero-shot: pool 331 ảnh, $\alpha=0$ từ MKG-W val. C: chỉ encoder XLM-R H/14; +KG: thêm kênh đếm
cạnh. Tất định. MRR nhân 100.}
\label{tab:vn-a0}
\footnotesize
\setlength{\tabcolsep}{2.5pt}
\begin{tabular}{@{}lrrrrrrr@{}}
\toprule
 & & \multicolumn{2}{c}{R@1} & \multicolumn{2}{c}{R@10} & \multicolumn{2}{c}{MRR} \\
\cmidrule(lr){3-4}\cmidrule(lr){5-6}\cmidrule(l){7-8}
Nhóm câu & $n$ & C & +KG & C & +KG & C & +KG \\
\midrule
Tất cả & 227 & 23,35 & \textbf{33,92} & 72,69 & \textbf{74,01} & 37,18 & \textbf{45,18} \\
Tự động & 197 & 21,32 & \textbf{29,95} & 70,05 & \textbf{72,59} & 35,21 & \textbf{41,61} \\
Tay, trực tiếp & 12 & \textbf{75,00} & 33,33 & \textbf{100,00} & 75,00 & \textbf{81,25} & 48,90 \\
Tay, gián tiếp & 18 & 11,11 & \textbf{77,78} & 83,33 & \textbf{88,89} & 29,41 & \textbf{81,81} \\
\bottomrule
\end{tabular}
\end{table}

\begin{table}[t]
\centering
\caption{A1: $\alpha$ chọn trên val tiếng Việt, test chạy một lần (103 câu tự động, pool 161 ảnh). Hàng cuối chỉ
để đối chiếu. MRR nhân 100.}
\label{tab:vn-a1}
\footnotesize
\begin{tabular}{@{}lrrrr@{}}
\toprule
Test & R@1 & R@5 & R@10 & MRR \\
\midrule
CLIP (XLM-R H/14) & 37,86 & 72,82 & 79,61 & 52,39 \\
+KG, $\alpha=0{,}7$ (chọn trên val) & 47,57 & 75,73 & 84,47 & 60,41 \\
+KG, $\alpha=0$ (của A0) & 50,49 & 70,87 & 81,55 & 60,46 \\
\bottomrule
\end{tabular}
\end{table}

Ở A0 (Bảng~\ref{tab:vn-a0}), kênh KG nâng R@1 từ 23,35 lên 33,92 (+10,6 điểm). Câu gián tiếp đúng ở hạng 1 tăng
từ 2 lên 14 trên 18; câu trực tiếp giảm từ 9 xuống 4 trên 12. Linker nối được ít nhất một thực thể ở 84,14\% số
câu.

Ở A1 (Bảng~\ref{tab:vn-a1}), R@1 trên val là 36,17 với $\alpha \in [0;\,0{,}35]$, 37,23 với
$\alpha \in [0{,}4;\,0{,}75]$ và 29,79 ở $\alpha=1$; hai đoạn phẳng chỉ cách nhau 1 câu trên 94, và MRR chọn
$\alpha=0{,}7$. Trên test, R@1 tăng từ 37,86 lên 47,57 (+9,7 điểm). Với $\alpha=0$, R@1 cao hơn (thêm 3 câu) nhưng R@5 và
R@10 thấp hơn, MRR gần bằng; $\alpha$ không được chọn lại theo test.

\subsection{Phân tích chuyển giao}\label{sec:vn-analysis}

\textbf{Kênh đếm cạnh áp dụng được trên dữ liệu mới} dù khác ngôn ngữ, khác encoder và đồ thị nhỏ hơn (325 so với
34\,196 triple KGC-train của MKG-W). Kênh không có tham số học, nên chỉ cần một linker theo tên, một KG và một $\alpha$; ở
A0, kênh này nâng R@1 thêm 10,6 điểm.

\textbf{Truy vấn gián tiếp.} Với câu ``Ngôi chùa ở Hà Nội được xây dựng dưới thời nhà Lý'', Chùa Một
Cột (đúng) lên hạng 1 nhờ ba cạnh \texttt{P131}~$\rightarrow$~Hà Nội, \texttt{P2348} (thời kỳ)~$\rightarrow$~nhà
Lý, \texttt{P31}~$\rightarrow$~chùa; vị trí hành chính và thời kỳ không thấy được trong ảnh. Bốn câu gián tiếp còn
sai đều liên quan tới quan hệ lớp con, mà kênh đếm cạnh trực tiếp (1-hop) không đi qua. Ở ba câu hỏi ``nhà thờ'' (Ninh Bình, Khánh Hòa, Thành phố Hồ
Chí Minh), từ ``nhà thờ'' được nối vào lớp Q16970, còn đáp án là nhà thờ chính tòa (lớp con Q2977), nên đáp án
thiếu cạnh \texttt{P31} này và không được cộng điểm hơn một nhà thờ thường cùng tỉnh. Câu còn lại nhắc
``di sản thế giới'' (Q9259), còn danh hiệu của đáp án là một thực thể khác (Q43113623).

\textbf{Truy vấn trực tiếp ở $\alpha=0$.} Linker nối từ chỉ loại (``nhà thờ'', ``tháp'',
``chùa'') vào thực thể lớp, nên ảnh có \texttt{P31} tới đúng lớp đó được cộng một cạnh, còn đáp án thường không có
(nhà thờ chính tòa thuộc lớp con; chùa có tháp không thuộc lớp ``tháp''). Với $\alpha=0$, chênh lệch một cạnh quyết
định thứ hạng, bất kể độ tương đồng CLIP với các chi tiết như ``sơn màu hồng'' hay ``hai tháp chuông''. Từ dùng để so sánh cũng bị nối: câu ``Công trình bằng
đá xám có nhiều tầng mái cong như đình chùa, trên đỉnh có cây thánh giá'' đưa Chùa Bút Tháp lên hạng 1 nhờ cạnh
\texttt{P31}~$\rightarrow$~chùa, trong khi đáp án là một nhà thờ chính tòa. Tỉ lệ câu nối được thực thể (84,14\%)
tương đương MKG-W test (80,94\%); các lỗi phân tích ở trên đều do liên kết từ chỉ loại vào sai lớp.

\textbf{$\alpha$ theo loại câu.} Trên val A1 (chỉ có câu tự động), các giá trị $\alpha$ từ 0 đến 0{,}75 chênh nhau 1
câu. Ở cùng $\alpha=0$, nhóm gián tiếp thêm 12 câu đúng, nhóm trực tiếp mất 5 câu; một $\alpha$ chung cho mọi câu không
tối ưu cho cả hai nhóm.

\subsection{Hạn chế}\label{sec:vn-limits}

\begin{itemize}
\item Cỡ mẫu nhỏ, chưa có khoảng tin cậy bootstrap; câu soạn tay chỉ được chấm ở $\alpha=0$.
\item Thiên lệch dữ liệu: ưu tiên di tích nổi tiếng; tỉnh lẫn trước và sau đợt sáp nhập năm 2025; thời kỳ chỉ suy
  theo năm.
\item Linker không chuẩn hoá hai kiểu đặt dấu thanh (``hòa''/``hoà''), dù cả hai kiểu đều có trong nhãn và câu;
  ảnh hưởng chưa đo. Chưa thử đi thêm \texttt{P279}, đếm 2-hop hay $\alpha$ theo từng câu.
\end{itemize}
